% ---------------------------------------------------------------
\section{Discussion}
\label{sec:discussion}

\subsection{Why the Framework Behaves as It Does}
The results follow from one identifiability fact. Unlabelled inputs constrain $P(X)$,
and they constrain accuracy only through assumptions that tie $P(Y\mid X)$ to the
source law. Calibrated average confidence makes this assumption explicit. Its error
stays small while the source correctness map holds (C0--C2 and C6), and it fails
silently when the map breaks, either because the label law changes (C3--C5) or because
the sensors drift beyond the calibration data (UCI). The input channel fails in the
opposite way: it alarms on harmless change (C1, seasonality) and saturates on natural
streams (every ELEC2 and NOAA batch). Keeping the two channels separate and scoring
both against verified harm exposes these complementary blind spots. Released labels
are the only evidence that resolves them, and the delay sets how soon they arrive.

\subsection{Practical Implications and Strengths}
Four recommendations follow. (i) An input alarm should not be presented as evidence
of harm, and the absence of an alarm should not be read as safety; the view should
state ``performance unknown'' until labels are released. (ii) Some labels should be
acquired at random with known inclusion probabilities, which removes the bias of
confidence-targeted labelling (Appendix~\ref{app:extra}). (iii) Released labels should
recalibrate the correctness map, and the benefit should be monitored rather than
assumed, since it held on UCI but not on ELEC2. (iv) Detectors should be tuned for
their stationary alarm burden as well as their sensitivity. The study's main strength
is that every channel is scored against verified harm under an enforced release
clock, with replication at 1,500 seeds, an exact clean rerun, and an internal baseline
for every component.

\subsection{Limitations and Failure Cases}
Most evidence comes from designed laws with eight features and two learners; more
seeds narrow uncertainty within these laws but do not make them representative of
other domains. Three natural streams with one split protocol each cannot support
general transfer claims. ELEC2's order is approximate, and UCI's release schedule is
simulated at the level of whole cohorts. The Poisson-binomial interval treats the
calibrated probabilities as known and independent; it covered only 9/21 NOAA batches,
23/55 ELEC2 batches and 0/9 UCI cohorts, so it is not a confidence guarantee. Feature
rankings and correlation pairs describe change and do not identify its cause, and the
aligned view was not evaluated with users. The documented failure cases are
conditional-only change, calibration transfer under sensor drift, ineffective
recalibration on ELEC2, alarm saturation on natural streams, and a concern score that
ranks future harm worse than chance on a held-out law (AUC 0.465).

\subsection{Threats to Validity}
\emph{Internal:} several extension hypotheses were formed after the first synthetic
results, so those extensions are exploratory, and several studies share seed banks
and cached predictions. \emph{Construct:} the harm rule uses a fixed margin
$\delta=0.05$ against a model-specific reference, so harm counts are not comparable
across learners. \emph{External:} the tabular, low-dimensional setting differs from
the deep setting of DriftLens and the high-dimensional, correlated setting of DDE.
\emph{Statistical conclusion:} intervals come from seed bootstraps; batches within a
stream are serially dependent and are never treated as independent units.

\subsection{Were the Research Gap and Objectives Met?}
Table~\ref{tab:rq} maps the Phase~2 research questions to the evidence. The
\emph{empirical} gap was met. The combination of input evidence, feature evidence and
calibrated impact was evaluated against verified harm under enforced delays, and its
failure modes were quantified. The \emph{decision-support} part of the gap was not met,
because no analyst study was run.

\begin{table}[t]
\centering
\caption{Status of the Phase~2 research questions.}
\label{tab:rq}
\setlength{\tabcolsep}{3pt}
\footnotesize
\begin{tabular}{L{3.2cm}L{1.15cm}L{3.6cm}}
\toprule
Research question & Status & Evidence and boundary \\
\midrule
RQ1: Does calibrated impact anticipate harm better than drift scores under delay? & partial & yes when $P(Y|X)$ is stable; no under conditional-only change or sensor drift; forecasting gains only on known laws \\
RQ2: Does combined evidence localize changes better? & partial & marginal features P@1 1.0 (C1) and 0.988 (C2); dependence pair 30/30; cluster regions near chance \\
RQ3: Do coordinated views improve analyst decisions? & not tested & view implemented and leak-checked; no users \\
RQ4: Do findings transfer across regimes and streams? & partial & two synthetic families, three natural streams; transfer not established \\
\bottomrule
\end{tabular}
\end{table}
